\scanpage{203}
\mspage{59}
\deleted{Si $u$ est $W_0$-lisse, il}\inserted{Dans les trois cas, on sait que $u$} est une $W_0$-équiv. univ. sss ses fibres sont $W_0$-connexes. Or la $W_0$-lissité implique que \deleted{\unsure{$\forall i$ de}} $\forall i\in\operatorname{Ob}I$, l'inclusion

\[
 I'_i\longrightarrow{}_{i\backslash}I'
\]

est une $W_0$-équivalence, d'autre part $\,{}_{i\backslash}I'\longrightarrow{}_{i\backslash}I$ étant $W_0$-asph. donc une $W_0$-équivalence\unsure{,} par hypothèse, $\,{}_{i\backslash}I'$ est $0$-connexe, donc $I'_i$ aussi. Le cas $u$ $W_0$-propre est dual. Dans le cas $u$ une $W_0$-fibration faible, on sait encore que \deleted{c'est une $W_0$-} $I'_i\mathbin{\unsure{-|\to}}{}_{i\backslash}I'$ est une $W_0$-équivalence, et on conclut comme dans le cas lisse (ou, au choix, via $I'_i\longrightarrow I'_{/i}$).

Mais l'idée me vient que \deleted{\unsure{des}}\inserted{les} foncteurs $W_0$-connexes pourraient être \inserted{\unsure{exactement}} \deleted{des} foncteurs de localisation ? Il revient au $=$ de dire qu'un foncteur $u\colon I'\mathbin{-}I$ \deleted{tel que} \uline{conservatif} (i.e. $u(\alpha)$ iso $\Rightarrow\alpha$ iso) qui est $W_0$-connexe, est une équivalence. Mais c'est faux, comme on voit en prenant $I=e$, $I'$ un groupoïde $0$-connexe et non $1$-connexe : $I'\mathbin{-}e$ est $W_0$-connexe, (et $=$ une $W_0$-équiv. universelle), mais ce n'est pas une équivalence, (et pas un foncteur de localisation).

\figure{\begin{tikzpicture}
 \draw[thick] (0,.38) .. controls (.12,.24) and (.30,.07) .. (.42,0);
\end{tikzpicture}}
% archive sheet number: 203
